\documentclass[11pt,a4paper]{article}

\usepackage[a4paper,margin=1in]{geometry}
\usepackage{amsmath}
\usepackage{amssymb}
\usepackage{array}
\usepackage{booktabs}
\usepackage{enumitem}
\usepackage{tabularx}
\usepackage{tikz}
\usepackage[T1]{fontenc}
\usepackage{lmodern}
\usepackage{xcolor}

\definecolor{bookletblue}{RGB}{30,70,120}
\definecolor{bookletgrey}{RGB}{90,90,90}

\newcommand{\markalloc}[1]{\nobreak\hfill\mbox{(#1 Marks)}}
\newcommand{\questiontotal}[1]{\begin{flushright}\textbf{Total #1 Marks}\end{flushright}}

\setlength{\parindent}{0pt}
\setlength{\parskip}{0.45em}
\setlist[enumerate,1]{label=(\alph*), leftmargin=1.1cm, itemsep=0.55em, topsep=0.4em}

\begin{document}

\begin{center}
{\Large\bfseries\color{bookletblue} Full Revision Booklet}\\[0.4em]
{\large 5CM522 Operational Research}\\[0.3em]
{\color{bookletgrey} Combined and renumbered from Practice Papers A and B}
\end{center}

\vspace{0.8em}
\hrule
\vspace{0.8em}

This booklet contains all questions from both current practice papers.
The questions have been renumbered so that the set reads as one continuous revision pack.

\begin{center}
\begin{tabular}{lll}
\toprule
Booklet Question & Source & Original Question \\
\midrule
1--6 & Practice Paper A & Questions 1--6 \\
7--12 & Practice Paper B & Questions 1--6 \\
\bottomrule
\end{tabular}
\end{center}

\newpage
\textbf{1.}

\textit{Source: Practice Paper A, original Question 1.}



Elmbridge Furniture makes two products each week: desks and shelving units.

Each desk uses 5 carpentry hours and 2 finishing hours.
Each shelving unit uses 3 carpentry hours and 4 finishing hours.
At most 120 carpentry hours and 96 finishing hours are available each week.
At least 8 shelving units must be produced to satisfy an online retailer.

The weekly contribution is \pounds58 per desk and \pounds46 per shelving unit.

\begin{enumerate}
\item Formulate a linear programming model for this problem, clearly defining the decision variables, objective function and constraints. \markalloc{6}

\item Solve the model to determine the weekly production plan that maximises contribution. You may use a graphical method, algebra or Python, but your work must be clear. \markalloc{6}

\item Identify the binding constraints and determine the slack in each constraint at the optimum. \markalloc{4}

\item The contribution on shelving units rises to \pounds52. Determine whether the optimal basis changes and justify your answer clearly. \markalloc{5}

\item A new agreement requires shelving units to make up at least 30\% of total weekly output. Reformulate the model and determine whether the original optimum remains feasible. If not, find the revised optimal plan. \markalloc{4}
\end{enumerate}

\questiontotal{25}

\newpage

\textbf{2.}

\textit{Source: Practice Paper A, original Question 2.}



GreenGrid Energy is choosing a set of improvement projects for a regional depot network.
Each project is either undertaken in full or not undertaken at all.

\begin{center}
\begin{tabular}{lcc}
\toprule
Project & Cost (\pounds m) & Annual Benefit Score \\
\midrule
$A$ & 6 & 14 \\
$B$ & 8 & 17 \\
$C$ & 5 & 11 \\
$D$ & 9 & 20 \\
$E$ & 4 & 9 \\
$F$ & 7 & 15 \\
\bottomrule
\end{tabular}
\end{center}

The total budget is \pounds21 million.

The following additional conditions apply:
\begin{itemize}[leftmargin=0.9cm, itemsep=0.25em]
\item if project $D$ is selected, project $B$ must also be selected;
\item projects $A$ and $F$ cannot both be selected;
\item at least one of projects $C$ and $E$ must be selected.
\end{itemize}

\begin{enumerate}
\item Formulate a binary integer programming model to maximise total annual benefit subject to the budget and project-selection rules. Clearly define all decision variables. \markalloc{8}

\item Solve the model in Python or by systematic analysis, and determine which projects should be selected and the maximum total annual benefit score. \markalloc{7}

\item Explain why this problem requires binary decision variables rather than the continuous variables used in ordinary linear programming. \markalloc{3}

\item Modify the model so that exactly four projects must be selected, then resolve or explain how the optimal solution would change. \markalloc{4}

\item Suggest one realistic extension that would make this capital-selection model more useful in practice. No full formulation is required. \markalloc{3}
\end{enumerate}

\questiontotal{25}

\newpage

\textbf{3.}

\textit{Source: Practice Paper A, original Question 3.}



The network below represents proposed communication links between seven depots.
The number on each edge is the installation cost in thousands of pounds.

\begin{center}
\begin{tabular}{lclc}
\toprule
Link & Cost & Link & Cost \\
\midrule
$A$--$B$ & 6 & $C$--$F$ & 6 \\
$A$--$C$ & 4 & $D$--$E$ & 5 \\
$A$--$D$ & 9 & $D$--$F$ & 7 \\
$B$--$C$ & 5 & $D$--$G$ & 8 \\
$B$--$E$ & 7 & $E$--$F$ & 3 \\
$C$--$D$ & 6 & $E$--$G$ & 4 \\
$C$--$E$ & 8 & $F$--$G$ & 5 \\
\bottomrule
\end{tabular}
\end{center}

\begin{enumerate}
\item Explain what is meant by a minimum spanning tree and why it is useful in network design. \markalloc{4}

\item Use either Prim's algorithm or Kruskal's algorithm to determine a minimum spanning tree for this network. Show the main steps of the algorithm. \markalloc{10}

\item State the selected links and the total installation cost of the spanning tree. \markalloc{4}

\item The firm now requires one additional backup link so that there is exactly one cycle in the final design. Identify the cheapest such addition and state the revised total cost. \markalloc{4}

\item Comment briefly on one limitation of using only an MST for a real logistics or telecoms network. \markalloc{3}
\end{enumerate}

\questiontotal{25}

\newpage

\textbf{4.}

\textit{Source: Practice Paper A, original Question 4.}



Westmoor Cars is deciding how many display bays to provide for newly delivered vehicles.
Each week, the number of cars delivered and the number of cars requested by customers are uncertain.

\begin{center}
\begin{tabular}{lccccc}
\toprule
Cars Delivered & 6 & 7 & 8 & 9 & 10 \\
\midrule
Probability & 0.10 & 0.20 & 0.35 & 0.25 & 0.10 \\
\bottomrule
\end{tabular}
\end{center}

\begin{center}
\begin{tabular}{lccccc}
\toprule
Weekly Demand & 5 & 6 & 7 & 8 & 9 \\
\midrule
Probability & 0.15 & 0.25 & 0.30 & 0.20 & 0.10 \\
\bottomrule
\end{tabular}
\end{center}

Each permanent display bay costs \pounds18 per week.
Any cars that cannot be placed in a display bay must be held in overflow storage at a cost of \pounds30 per car per week.

For consistency, use the following random numbers for a 12-week simulation:

\begin{center}
\scriptsize
\begin{tabular}{lcccccccccccc}
\toprule
Week & 1 & 2 & 3 & 4 & 5 & 6 & 7 & 8 & 9 & 10 & 11 & 12 \\
\midrule
Delivery RN & 0.119 & 0.503 & 0.512 & 0.860 & 0.103 & 0.223 & 0.601 & 0.557 & 0.783 & 0.548 & 0.731 & 0.768 \\
Demand RN & 0.751 & 0.587 & 0.240 & 0.614 & 0.111 & 0.817 & 0.450 & 0.815 & 0.685 & 0.679 & 0.210 & 0.252 \\
\bottomrule
\end{tabular}
\normalsize
\end{center}

\begin{enumerate}
\item Describe how you would model this problem as a simulation study. Your answer should identify the stochastic inputs, state variables and performance measures. \markalloc{5}

\item Simulate 12 weeks of operation for a capacity of 7 display bays using the random numbers above. You may use Excel or Python to carry out the calculations, but your logic must be clear. \markalloc{8}

\item Estimate the average weekly cost for a capacity of 7 bays. \markalloc{3}

\item Repeat the analysis for capacities of 6, 8 and 9 bays using the same random numbers, then recommend a capacity. Present your results in a clearly labelled table. \markalloc{6}

\item State two reasons why a much longer simulation or repeated runs would improve confidence in the recommendation. \markalloc{3}
\end{enumerate}

\questiontotal{25}

\newpage

\textbf{5.}

\textit{Source: Practice Paper A, original Question 5.}



The demand analyst at Calder Kitchens has collected the following quarterly sales data for one of its product ranges.

\begin{center}
\begin{tabular}{cccc}
\toprule
Period & Year & Quarter & Demand \\
\midrule
1  & 2020 & Q1 & 510 \\
2  & 2020 & Q2 & 620 \\
3  & 2020 & Q3 & 570 \\
4  & 2020 & Q4 & 690 \\
5  & 2021 & Q1 & 548 \\
6  & 2021 & Q2 & 662 \\
7  & 2021 & Q3 & 608 \\
8  & 2021 & Q4 & 742 \\
9  & 2022 & Q1 & 590 \\
10 & 2022 & Q2 & 708 \\
11 & 2022 & Q3 & 646 \\
12 & 2022 & Q4 & 792 \\
13 & 2023 & Q1 & 636 \\
14 & 2023 & Q2 & 760 \\
15 & 2023 & Q3 & 692 \\
16 & 2023 & Q4 & 852 \\
\bottomrule
\end{tabular}
\end{center}

\begin{enumerate}
\item Using periods 1--12 only, calculate suitable centred 4-quarter moving averages and estimate the trend where possible. Plot the original data together with the estimated trend. \markalloc{8}

\item Using periods 1--12 only, calculate seasonal indices using both additive and multiplicative decomposition. \markalloc{7}

\item Use periods 13--16 as a validation set. Produce forecasts for those periods using both models and compare their accuracy using one suitable error measure. \markalloc{6}

\item Using the better model, forecast the four quarters of 2024 and justify your model choice briefly. \markalloc{4}
\end{enumerate}

\questiontotal{25}

\newpage

\textbf{6.}

\textit{Source: Practice Paper A, original Question 6.}



Northfield Relief ships medical supplies from three depots to four clinics through two regional hubs.
Supplies may travel only from a depot to a hub and then from a hub to a clinic.

\begin{center}
\begin{tabular}{lccc}
\toprule
Depot & Supply & Cost to $H_1$ & Cost to $H_2$ \\
\midrule
$D_1$ & 100 & 5 & 7 \\
$D_2$ & 90 & 4 & 6 \\
$D_3$ & 110 & 6 & 3 \\
\bottomrule
\end{tabular}
\end{center}

\begin{center}
\begin{tabular}{lccccc}
\toprule
Hub & Cost to $C_1$ & Cost to $C_2$ & Cost to $C_3$ & Cost to $C_4$ & Total Demand \\
\midrule
$H_1$ & 4 & 6 & 5 & 7 & \\
$H_2$ & 6 & 4 & 3 & 5 & \\
\midrule
Demand & 60 & 70 & 80 & 75 & 285 \\
\bottomrule
\end{tabular}
\end{center}

\begin{enumerate}
\item Formulate a transhipment model for this problem. Clearly define the decision variables, objective function, supply constraints, demand constraints and hub flow-balance constraints. \markalloc{8}

\item Solve the model and determine the optimal shipment plan, the total transportation cost and any unused supply. Report only the non-zero shipments. \markalloc{7}

\item A disruption means that hub $H_1$ can process at most 130 units, and the route from $H_2$ to clinic $C_3$ is unavailable. Show how the model changes and determine the revised total cost. State only the revised non-zero shipments or the flows that change. \markalloc{6}

\item In the disrupted scenario, up to 20 extra units of processing capacity can be bought for hub $H_1$ at a one-off cost of \pounds50. Determine whether this is worthwhile and justify your answer. \markalloc{4}
\end{enumerate}

\questiontotal{25}

\newpage

\textbf{7.}

\textit{Source: Practice Paper B, original Question 1.}



Silverforge Engineering is studying two scheduling problems.

\textbf{Part A: Flow Shop Scheduling}

All five jobs must be processed first on Machine 1 and then on Machine 2.

\begin{center}
\begin{tabular}{lcc}
\toprule
Job & Machine 1 Time & Machine 2 Time \\
\midrule
$J_1$ & 7 & 5 \\
$J_2$ & 3 & 8 \\
$J_3$ & 6 & 4 \\
$J_4$ & 2 & 9 \\
$J_5$ & 5 & 6 \\
\bottomrule
\end{tabular}
\end{center}

\textbf{Part B: Job Shop Scheduling}

Each job has its own machine order and must follow that order exactly.

\begin{center}
\begin{tabular}{lcc}
\toprule
Job & Sequence on Machines & Processing Times \\
\midrule
$A$ & $M_1, M_3, M_2$ & $4, 5, 3$ \\
$B$ & $M_2, M_1, M_3$ & $3, 4, 6$ \\
$C$ & $M_3, M_2, M_1$ & $5, 2, 4$ \\
$D$ & $M_1, M_2, M_3$ & $2, 6, 3$ \\
\bottomrule
\end{tabular}
\end{center}

\begin{enumerate}
\item Use Johnson's Rule to determine an optimal production sequence for the flow shop problem in Part A. Show the main steps. \markalloc{6}

\item Draw a Gantt chart for your sequence and calculate the makespan. \markalloc{5}

\item Compare your Johnson schedule with the original order $J_1,J_2,J_3,J_4,J_5$ and comment on the improvement. \markalloc{4}

\item Solve the job shop scheduling problem in Part B and state an optimal schedule and the corresponding makespan. You may present the result as a table of operation start and finish times or as a Gantt chart. \markalloc{6}

\item Suppose job $C$ must finish by time 15. Explain how the model changes and discuss the possible effect on feasibility and makespan. \markalloc{4}
\end{enumerate}

\questiontotal{25}

\newpage

\textbf{8.}

\textit{Source: Practice Paper B, original Question 2.}



County Response wants to choose a limited number of patrol bases.
The aim is not to cover every district, but to maximise the total population covered while opening at most three bases.

\begin{center}
\begin{tabular}{lccc}
\toprule
Base & Cost & Population Covered & Districts Covered \\
\midrule
$A$ & 14 & 18{,}000 & $\{1,2,4\}$ \\
$B$ & 11 & 15{,}000 & $\{1,3\}$ \\
$C$ & 13 & 21{,}000 & $\{2,3,5\}$ \\
$D$ & 10 & 12{,}000 & $\{4,5\}$ \\
$E$ & 12 & 16{,}000 & $\{2,5,6\}$ \\
$F$ & 9  & 14{,}000 & $\{3,4,6\}$ \\
\bottomrule
\end{tabular}
\end{center}

The population living in each district is given below.

\begin{center}
\begin{tabular}{lcccccc}
\toprule
District & 1 & 2 & 3 & 4 & 5 & 6 \\
\midrule
Population & 9{,}000 & 12{,}000 & 14{,}000 & 10{,}000 & 11{,}000 & 8{,}000 \\
\bottomrule
\end{tabular}
\end{center}

The budget is \pounds31 and at most 3 bases may be opened.

\begin{enumerate}
\item Formulate a binary integer programming model to maximise covered population subject to the budget and base limit. Clearly define any coverage variables you use. \markalloc{9}

\item Solve the model and identify which bases should be opened and the resulting total covered population. \markalloc{6}

\item Explain whether some districts may remain uncovered in an optimal solution and why this is acceptable in this model but not in a set covering model. \markalloc{4}

\item Modify the model so that district 6 must be covered and bases $B$ and $F$ cannot both be chosen. State the revised model changes clearly. \markalloc{6}
\end{enumerate}

\questiontotal{25}

\newpage

\textbf{9.}

\textit{Source: Practice Paper B, original Question 3.}



A university library help desk is modelled as a queueing system.
Students arrive randomly and service times are assumed to follow an exponential distribution.
You should adapt the SimPy structures used in class, but tailor the model to this setting.

Use the following assumptions:

\begin{center}
\begin{tabular}{lc}
\toprule
Quantity & Value \\
\midrule
Simulation length & 360 minutes \\
Average service time & 7 minutes \\
Baseline arrival rate & 10 students per hour \\
Peak arrival rate & 15 students per hour \\
Service target & 97\% wait less than 8 minutes \\
Priority users & 10\% of arrivals \\
Suggested random seed & 24 \\
\bottomrule
\end{tabular}
\end{center}

\begin{enumerate}
\item Describe the main elements of a discrete-event simulation model for this system: entities, resources, events, stochastic inputs and performance measures. \markalloc{5}

\item Adapt a SimPy model to represent the baseline system with 2 advisers and an arrival rate of 10 students per hour. Report the total served, average wait, maximum wait and percentage waiting less than 8 minutes. \markalloc{6}

\item Investigate the peak-period scenario with 2, 3 and 4 advisers. Present your results in a clearly labelled table and determine the minimum number of advisers needed to meet the service target. \markalloc{4}

\item Extend the model so that 10\% of users are priority users who should be served ahead of standard users. Explain how the model should be changed and discuss the likely effect on standard-user waits. \markalloc{6}

\item Explain why a single run is not enough to support a staffing decision and describe how you would use replications and/or a warm-up period to improve the analysis. \markalloc{4}
\end{enumerate}

\questiontotal{25}

\newpage

\textbf{10.}

\textit{Source: Practice Paper B, original Question 4.}



The operations analyst at Broadleaf Homeware has collected the following quarterly demand data.

\begin{center}
\begin{tabular}{cccc}
\toprule
Period & Year & Quarter & Demand \\
\midrule
1  & 2019 & Q1 & 420 \\
2  & 2019 & Q2 & 500 \\
3  & 2019 & Q3 & 470 \\
4  & 2019 & Q4 & 560 \\
5  & 2020 & Q1 & 450 \\
6  & 2020 & Q2 & 536 \\
7  & 2020 & Q3 & 500 \\
8  & 2020 & Q4 & 600 \\
9  & 2021 & Q1 & 485 \\
10 & 2021 & Q2 & 575 \\
11 & 2021 & Q3 & 535 \\
12 & 2021 & Q4 & 645 \\
13 & 2022 & Q1 & 522 \\
14 & 2022 & Q2 & 618 \\
15 & 2022 & Q3 & 572 \\
16 & 2022 & Q4 & 694 \\
\bottomrule
\end{tabular}
\end{center}

\begin{enumerate}
\item Using periods 1--12 only, calculate suitable centred 4-quarter moving averages and estimate the trend values where possible. Plot the original data together with the estimated trend. \markalloc{8}

\item Using periods 1--12 only, calculate seasonal indices using both additive and multiplicative decomposition models. \markalloc{7}

\item Use the 2022 data (periods 13--16) as a validation set. Generate forecasts for these four periods using both models and compare their accuracy using one suitable measure such as MAD, MSE or MAPE. \markalloc{6}

\item Using the better model, produce forecasts for the four quarters of 2023 and justify your final model choice. \markalloc{4}
\end{enumerate}

\questiontotal{25}

\newpage

\textbf{11.}

\textit{Source: Practice Paper B, original Question 5.}



Harbour Plastics makes two products each week: casings and inserts.

Each casing requires 4 moulding hours and 1 assembly hour.
Each insert requires 2 moulding hours and 3 assembly hours.
At most 100 moulding hours and 90 assembly hours are available.
No more than 24 inserts can be sold each week.

The weekly contribution is \pounds40 per casing and \pounds34 per insert.

\begin{enumerate}
\item Formulate a linear programming model for this problem. \markalloc{6}

\item Solve the model to determine the production plan that maximises weekly contribution. \markalloc{6}

\item State the binding constraints and the unused capacity in each department. \markalloc{4}

\item A distributor now requires at least 25\% of total output to be inserts. Reformulate the model and determine whether the original solution remains feasible. \markalloc{5}

\item If up to 12 extra assembly hours can be purchased at \pounds3 per hour for one week only, determine whether they should be bought. \markalloc{4}
\end{enumerate}

\questiontotal{25}

\newpage

\textbf{12.}

\textit{Source: Practice Paper B, original Question 6.}



Northline Distribution must ship goods from three warehouses to four stores at minimum cost.

\begin{center}
\begin{tabular}{lcccc}
\toprule
Warehouse & $S_1$ & $S_2$ & $S_3$ & $S_4$ \\
\midrule
$W_1$ & 6 & 8 & 10 & 7 \\
$W_2$ & 9 & 5 & 7 & 6 \\
$W_3$ & 4 & 7 & 6 & 8 \\
\bottomrule
\end{tabular}
\end{center}

\begin{center}
\begin{tabular}{lc}
\toprule
Warehouse & Supply \\
\midrule
$W_1$ & 80 \\
$W_2$ & 100 \\
$W_3$ & 90 \\
\bottomrule
\end{tabular}
\hspace{1.5cm}
\begin{tabular}{lc}
\toprule
Store & Demand \\
\midrule
$S_1$ & 70 \\
$S_2$ & 60 \\
$S_3$ & 80 \\
$S_4$ & 60 \\
\bottomrule
\end{tabular}
\end{center}

\begin{enumerate}
\item Formulate a transportation model for this problem, clearly defining the decision variables, objective function and constraints. \markalloc{7}

\item Solve the model and determine the optimal shipping plan and total transportation cost. Report only the non-zero shipments. \markalloc{7}

\item Resolve the model if warehouse $W_2$ cannot supply store $S_2$, and at least 20 units must be shipped from $W_1$ to $S_4$. Comment on the impact on total cost. \markalloc{6}

\item Explain one realistic extension that would make this transportation model more useful in practice. \markalloc{5}
\end{enumerate}

\questiontotal{25}
\end{document}
